\chapter{2007年浙江卷（文）}

\FigureLayoutDeclare{img/2007/zhejiang_liberal/q5_fig1.png}{5.0cm}{54.24bp 2.64bp 60.96bp 1.92bp}

\section{单选题}

\begin{problem}
设全集 \(U=\{1, 3, 5, 6, 8\}\)，\(A=\{1, 6\}\)，\(B=\{5, 6, 8\}\)，则 \(\left(\complement_U A\right)\cap B=\)（\quad）

\choices
    {\(\{6\}\)}
    {\(\{5, 8\}\)}
    {\(\{6, 8\}\)}
    {\(\{3, 5, 6, 8\}\)}
\end{problem}
\answer{B}

\begin{solution}
由全集和集合 \(A\) 可得 \(\complement_U A=\{3,5,8\}\). 因此 \(\left(\complement_U A\right)\cap B=\{3,5,8\}\cap\{5,6,8\}=\{5,8\}\)，故选 B.
\end{solution}

\begin{problem}
已知 \(\cos\left(\frac{\pi}{2}+\varphi\right)=\frac{\sqrt{3}}{2}\)，且 \(\left|\varphi\right|<\frac{\pi}{2}\)，则 \(\tan\varphi=\)（\quad）

\choices
    {\(-\frac{\sqrt{3}}{3}\)}
    {\(\frac{\sqrt{3}}{3}\)}
    {\(-\sqrt{3}\)}
    {\(\sqrt{3}\)}
\end{problem}
\answer{C}

\begin{solution}
由 \(\cos\left(\frac{\pi}{2}+\varphi\right)=-\sin\varphi\)，得 \(\sin\varphi=-\frac{\sqrt{3}}{2}\). 又因 \(\left|\varphi\right|<\frac{\pi}{2}\)，所以 \(\cos\varphi>0\)，从而 \(\cos\varphi=\frac{1}{2}\). 因此 \(\tan\varphi=\frac{\sin\varphi}{\cos\varphi}=-\sqrt{3}\)，故选 C.
\end{solution}

\begin{problem}
“\(x>1\)” 是 “\(x^2>x\)” 的（\quad）

\choices
    {充分而不必要条件}
    {必要而不充分条件}
    {充分必要条件}
    {既不充分也不必要条件}
\end{problem}
\answer{A}

\begin{solution}
若 \(x>1\)，则 \(x(x-1)>0\)，即 \(x^2>x\)，所以 \(x>1\) 是 \(x^2>x\) 的充分条件. 反之，\(x^2>x\) 等价于 \(x(x-1)>0\)，解得 \(x<0\) 或 \(x>1\). 例如 \(x=-1\) 时 \(x^2>x\) 成立而 \(x>1\) 不成立，因此不是必要条件，故选 A.
\end{solution}

\begin{problem}
直线 \(x - 2y + 1 = 0\) 关于直线 \(x = 1\) 对称的直线方程是（\quad）

\choices
    {\(x + 2y - 1 = 0\)}
    {\(2x + y - 1 = 0\)}
    {\(2x + y - 3 = 0\)}
    {\(x + 2y - 3 = 0\)}
\end{problem}
\answer{D}

\begin{solution}
设原直线上的点为 \((x,y)\)，关于直线 \(x=1\) 对称后的点为 \((X,Y)\). 由轴对称关系，\(X=2-x\)，\(Y=y\)，即 \(x=2-X\)，\(y=Y\). 代入原直线方程得 \((2-X)-2Y+1=0\)，化简为 \(X+2Y-3=0\). 将 \(X\)，\(Y\) 仍记为 \(x\)，\(y\)，所求直线为 \(x+2y-3=0\)，故选 D.
\end{solution}

\begin{problem}
要在边长为 \(16\text{ 米}\) 的正方形草坪上安装喷水龙头，使整个草坪都能喷洒到水. 假设每个喷水龙头的喷洒范围都是半径为 \(6\text{ 米}\) 的圆面，则需安装这种喷水龙头的个数最少是（\quad）

\bitmapfigure[width=0.44\linewidth]{img/2007/zhejiang_liberal/q5_fig1.png}

\choices
    {\(6\)}
    {\(5\)}
    {\(4\)}
    {\(3\)}
\end{problem}
\answer{C}

\begin{solution}
正方形的四个顶点中，任意两个不同顶点的距离至少为 \(16\)，而一个半径为 \(6\) 的圆的直径为 \(12\)，所以一个喷水范围不可能同时覆盖两个顶点. 四个顶点都必须被覆盖，故喷水龙头至少需要 \(4\) 个.

将正方形看作四个边长为 \(8\) 的小正方形，并在它们的中心分别安装喷水龙头. 每个小正方形内任一点到其中心的距离不超过 \(4\sqrt{2}<6\)，所以四个喷水范围能够覆盖整个草坪. 因而最少需要 \(4\) 个，故选 C.
\end{solution}

\begin{problem}
\(\left(\sqrt{x} - \frac{1}{x}\right)^9\) 展开式中的常数项是（\quad）

\choices
    {\(-36\)}
    {\(36\)}
    {\(-84\)}
    {\(84\)}
\end{problem}
\answer{C}

\begin{solution}
展开式的通项为 \(\mathrm C_9^k\left(\sqrt{x}\right)^{9-k}\left(-\frac{1}{x}\right)^k=(-1)^k\mathrm C_9^k x^{\frac{9-3k}{2}}\)，其中 \(k=0\)，\(1\)，\(\ldots\)，\(9\). 常数项要求 \(\frac{9-3k}{2}=0\)，所以 \(k=3\). 因而常数项为 \((-1)^3\mathrm C_9^3=-84\)，故选 C.
\end{solution}

\begin{problem}
若 \(P\) 是两条异面直线 \(l\)，\(m\) 外的任意一点，则（\quad）

\choices
    {过点 \(P\) 有且仅有一条直线与 \(l\)，\(m\) 都平行}
    {过点 \(P\) 有且仅有一条直线与 \(l\)，\(m\) 都垂直}
    {过点 \(P\) 有且仅有一条直线与 \(l\)，\(m\) 都相交}
    {过点 \(P\) 有且仅有一条直线与 \(l\)，\(m\) 都异面}
\end{problem}
\answer{B}

\begin{solution}
设异面直线 \(l\)，\(m\) 的方向向量分别为非共线向量 \(\bs{u}\)，\(\bs{v}\). 则 \(\bs{u}\times\bs{v}\) 同时垂直于 \(\bs{u}\) 和 \(\bs{v}\). 所有同时垂直于 \(\bs{u}\)，\(\bs{v}\) 的非零向量都与 \(\bs{u}\times\bs{v}\) 平行，因此过点 \(P\) 且方向为 \(\bs{u}\times\bs{v}\) 的直线唯一，并且与 \(l\)，\(m\) 都垂直，故选 B.
\end{solution}

\begin{problem}
甲，乙两人进行乒乓球比赛，比赛规则为“\(3\) 局 \(2\) 胜”，即以先赢 \(2\) 局者为胜. 根据经验，每局比赛中甲获胜的概率为 \(0.6\)，则本次比赛甲获胜的概率是（\quad）

\choices
    {\(0.216\)}
    {\(0.36\)}
    {\(0.432\)}
    {\(0.648\)}
\end{problem}
\answer{D}

\begin{solution}
甲获胜的比赛过程只有三种：前两局连胜，或前两局一胜一负且第三局获胜. 因而甲获胜的概率为 \(0.6^2+2\times0.6\times0.4\times0.6=0.36+0.288=0.648\)，故选 D.
\end{solution}

\begin{problem}
若非零向量 \(\bs{a}\)，\(\bs{b}\) 满足 \(\left|\bs{a} - \bs{b}\right| = \left|\bs{b}\right|\)，则（\quad）

\choices
    {\(\left|2\bs{b}\right| > \left|\bs{a} - 2\bs{b}\right|\)}
    {\(\left|2\bs{b}\right| < \left|\bs{a} - 2\bs{b}\right|\)}
    {\(\left|2\bs{a}\right| > \left|2\bs{a} - \bs{b}\right|\)}
    {\(\left|2\bs{a}\right| < \left|2\bs{a} - \bs{b}\right|\)}
\end{problem}
\answer{A}

\begin{solution}
将条件两边平方，得
\[
\left|\bs{a}\right|^2-2\bs{a}\cdot\bs{b}+\left|\bs{b}\right|^2=\left|\bs{b}\right|^2
\]
所以 \(2\bs{a}\cdot\bs{b}=\left|\bs{a}\right|^2\). 于是
\[
\left|2\bs{b}\right|^2-\left|\bs{a}-2\bs{b}\right|^2=4\left|\bs{b}\right|^2-\left(\left|\bs{a}\right|^2-4\bs{a}\cdot\bs{b}+4\left|\bs{b}\right|^2\right)=4\bs{a}\cdot\bs{b}-\left|\bs{a}\right|^2=\left|\bs{a}\right|^2>0
\]
因此 \(\left|2\bs{b}\right|>\left|\bs{a}-2\bs{b}\right|\)，故选 A.
\end{solution}

\begin{problem}
已知双曲线 \(\frac{x^2}{a^2} - \frac{y^2}{b^2} = 1\)（\(a>0\)，\(b>0\)）的左，右焦点分别为 \(F_1\)，\(F_2\)，\(P\) 是准线上一点，且 \(PF_1 \perp PF_2\)，\(\left|PF_1\right| \cdot \left|PF_2\right| = 4ab\)，则双曲线的离心率是（\quad）

\choices
    {\(\sqrt{2}\)}
    {\(\sqrt{3}\)}
    {\(2\)}
    {\(3\)}
\end{problem}
\answer{B}

\begin{solution}
设双曲线的焦距参数为 \(c\)，则 \(c^2=a^2+b^2\)，离心率为 \(e=\frac{c}{a}\). 由对称性，只需设 \(P\) 在右准线 \(x=\frac{a^2}{c}\) 上，记 \(P\left(\frac{a^2}{c},t\right)\). 令 \(F_1=(-c,0)\)，\(F_2=(c,0)\)，则
\(\overrightarrow{PF_1}=\left(-c-\frac{a^2}{c},-t\right)\)，\(\overrightarrow{PF_2}=\left(c-\frac{a^2}{c},-t\right)\).
由 \(PF_1\perp PF_2\)，得
\[
0=\overrightarrow{PF_1}\cdot\overrightarrow{PF_2}=-\left(c^2-\frac{a^4}{c^2}\right)+t^2
\]
所以 \(t^2=\frac{b^2(c^2+a^2)}{c^2}\). 代入两段距离，得
\(\left|PF_1\right|^2=2(c^2+a^2)\)，\(\left|PF_2\right|^2=2b^2\).
因此 \(\left|PF_1\right|\left|PF_2\right|=2b\sqrt{c^2+a^2}=4ab\)，从而 \(c^2+a^2=4a^2\). 于是 \(c^2=3a^2\)，故 \(e=\sqrt{3}\)，选 B.
\end{solution}

\section{填空题}

\begin{problem}
函数 \(y=\frac{x^2}{x^2+1}\)（\(x\in\R\)）的值域是\fillinblank{}.
\end{problem}
\answer{\([0,1)\)}

\begin{solution}
令 \(t=x^2\)，则 \(t\geqslant0\)，且 \(y=\frac{t}{t+1}\). 因为 \(t+1>0\)，所以 \(y\geqslant0\)，又有 \(y=1-\frac{1}{t+1}<1\). 反过来，对于任意 \(y_0\in[0,1)\)，取 \(t=\frac{y_0}{1-y_0}\geqslant0\)，再取 \(x=\sqrt{t}\)，就有 \(\frac{x^2}{x^2+1}=y_0\). 因此值域为 \([0,1)\).
\end{solution}

\begin{problem}
若 \(\sin \theta + \cos \theta = \frac{1}{5}\)，则 \(\sin 2\theta\) 的值是\fillinblank{}.
\end{problem}
\answer{\(-\frac{24}{25}\)}

\begin{solution}
两边平方，得 \(\sin^2\theta+2\sin\theta\cos\theta+\cos^2\theta=\frac{1}{25}\). 利用 \(\sin^2\theta+\cos^2\theta=1\) 及 \(2\sin\theta\cos\theta=\sin2\theta\)，可得 \(1+\sin2\theta=\frac{1}{25}\)，所以 \(\sin2\theta=-\frac{24}{25}\).
\end{solution}

\begin{problem}
某校有学生 \(2000\text{ 人}\)，其中高三学生 \(500\text{ 人}\). 为了解学生的身体素质情况，采用按年级分层抽样的方法，从该校学生中抽取一个 \(200\text{ 人}\) 的样本. 则样本中高三学生的人数为\fillinblank{}.
\end{problem}
\answer{\(50\)}

\begin{solution}
按年级分层并按各年级人数比例抽样，所以样本中高三学生人数为 \(200\times\frac{500}{2000}=50\).
\end{solution}

\begin{problem}
\(z = 2x + y\) 中的 \(x\)，\(y\) 满足约束条件 \(\begin{cases}
x - 2y + 5 \geqslant 0,\\
3 - x \geqslant 0,\\
x + y \geqslant 0,
\end{cases}\) 则 \(z\) 的最小值是\fillinblank{}.
\end{problem}
\answer{\(-\frac{5}{3}\)}

\begin{solution}
由 \(x+y\geqslant0\) 得 \(y\geqslant-x\)，所以 \(z=2x+y\geqslant x\). 又由 \(x-2y+5\geqslant0\) 和 \(y\geqslant-x\)，得 \(x-2y+5\leqslant x+2x+5=3x+5\)，从而 \(3x+5\geqslant0\)，即 \(x\geqslant-\frac{5}{3}\). 因此 \(z\geqslant-\frac{5}{3}\).

当 \(x=-\frac{5}{3}\)，\(y=\frac{5}{3}\) 时，三个约束分别成立，且 \(x+y=0\)，所以 \(z=2\left(-\frac{5}{3}\right)+\frac{5}{3}=-\frac{5}{3}\). 故最小值为 \(-\frac{5}{3}\).
\end{solution}

\begin{problem}
曲线 \( y = x^3 - 2x^2 - 4x + 2 \) 在点 \(\left(1,-3\right)\) 处的切线方程是\fillinblank{}.
\end{problem}
\answer{\(5x+y-2=0\)}

\begin{solution}
设 \(f(x)=x^3-2x^2-4x+2\)，则 \(f'(x)=3x^2-4x-4\)，所以 \(f'(1)=-5\). 切线经过 \((1,-3)\)，故方程为 \(y+3=-5(x-1)\)，即 \(5x+y-2=0\).
\end{solution}

\begin{problem}
某书店有 \(11\) 种杂志，\(2\text{ 元}\) \(1\text{ 本}\)的 \(8\) 种，\(1\text{ 元}\) \(1\text{ 本}\)的 \(3\) 种. 小张用 \(10\text{ 元}\)钱买杂志（每种至多买一本，\(10\text{ 元}\)钱刚好用完），则不同买法的种数是\fillinblank{}.（用数字作答）
\end{problem}
\answer{\(266\)}

\begin{solution}
设购买 \(2\) 元杂志 \(r\) 种，购买 \(1\) 元杂志 \(s\) 种，则 \(2r+s=10\)，且 \(0\leqslant r\leqslant8\)，\(0\leqslant s\leqslant3\). 因而只有 \((r,s)=(4,2)\) 或 \((5,0)\). 相应的买法数为
\[
\mathrm C_8^4\mathrm C_3^2+\mathrm C_8^5\mathrm C_3^0=70\times3+56=266
\]
所以不同买法的种数为 \(266\).
\end{solution}

\begin{problem}
已知点 \(O\) 在二面角 \( \alpha - AB - \beta \) 的棱上，点 \(P\) 在 \( \alpha \) 内，且 \(\angle POB = 45^{\circ}\). 若对于 \( \beta \) 内异于 \(O\) 的任意一点 \(Q\)，都有 \(\angle POQ \geqslant 45^{\circ}\)，则二面角 \( \alpha - AB - \beta \) 的大小是\fillinblank{}.
\end{problem}
\answer{\(90^{\circ}\)}

\begin{solution}
取 \(AB\) 为 \(x\) 轴，取 \(\alpha\) 为 \(xy\) 平面，并设二面角 \(\alpha-AB-\beta\) 的大小为 \(\omega\). 由 \(\angle POB=45^{\circ}\)，可在适当取向下写成 \(\overrightarrow{OP}=r\left(\frac{1}{\sqrt2},\frac{1}{\sqrt2},0\right)\)，其中 \(r>0\).

平面 \(\beta\) 内从点 \(O\) 出发的单位方向向量可写为 \(\left(\cos t,\sin t\cos\omega,\sin t\sin\omega\right)\). 因而 \(\overrightarrow{OP}\) 与该方向向量的内积为 \(\frac{r}{\sqrt2}(\cos t+\sin t\cos\omega)\). 当 \(t\) 变化时，其最大值为 \(\frac{r}{\sqrt2}\sqrt{1+\cos^2\omega}\)，所以 \(\beta\) 内射线与 \(OP\) 的最小夹角 \(\gamma\) 满足
\[
\cos\gamma=\frac{\sqrt{1+\cos^2\omega}}{\sqrt2}
\]
题设对任意 \(Q\) 都有 \(\angle POQ\geqslant45^{\circ}\)，即 \(\gamma\geqslant45^{\circ}\). 由于余弦函数在 \([0,\pi]\) 上单调递减，得 \(\frac{\sqrt{1+\cos^2\omega}}{\sqrt2}\leqslant\frac{1}{\sqrt2}\)，从而 \(\cos\omega=0\). 因此 \(\omega=90^{\circ}\).
\end{solution}

\section{解答题}

\begin{problem}
已知 \(\triangle ABC\) 的周长为 \(\sqrt{2} + 1\)，且 \(\sin A + \sin B = \sqrt{2} \sin C\).

\begin{enumerate}
\item 求边 \(AB\) 的长；
\item 若 \(\triangle ABC\) 的面积为 \(\frac{1}{6} \sin C\)，求角 \(C\) 的度数.
\end{enumerate}
\end{problem}

\begin{solution}
\begin{enumerate}
\item 设 \(BC=a\)，\(CA=b\)，\(AB=c\)，外接圆半径为 \(R\). 由正弦定理，\(a=2R\sin A\)，\(b=2R\sin B\)，\(c=2R\sin C\)，所以由题意得
\[
a+b=\sqrt{2}c
\]
又因为三角形的周长为 \(\sqrt{2}+1\)，所以
\[
(\sqrt{2}+1)c=a+b+c=\sqrt{2}+1
\]
从而 \(c=1\)，即 \(AB=1\).

\item 由三角形面积公式及 \(0<C<\pi\)，得
\[
\frac{1}{2}ab\sin C=\frac{1}{6}\sin C
\]
所以 \(ab=\frac{1}{3}\). 又由（1）知 \(a+b=\sqrt{2}\)，故
\[
a^2+b^2=(a+b)^2-2ab=2-\frac{2}{3}=\frac{4}{3}
\]
在 \(\triangle ABC\) 中，由余弦定理，
\[
\cos C=\frac{a^2+b^2-c^2}{2ab}=\frac{\frac{4}{3}-1}{\frac{2}{3}}=\frac{1}{2}
\]
因为 \(0<C<\pi\)，所以 \(C=60^{\circ}\).
\end{enumerate}
\end{solution}

\begin{problem}
已知数列 \(\{a_{n}\}\) 中的相邻两项 \(a_{2k-1}\)，\(a_{2k}\) 是关于 \(x\) 的方程 \(x^2 - (3k + 2^k)x + 3k \cdot 2^k = 0\) 的两个根，且 \(a_{2k-1} \leqslant a_{2k}\)（\(k=1\)，\(2\)，\(3\)，\(\cdots\)）.

\begin{enumerate}
\item 求 \(a_{1}\)，\(a_{3}\)，\(a_{5}\)，\(a_{7}\) 及 \(a_{2n}\)（\(n \geqslant 4\)）（不必证明）；
\item 求数列 \(\{a_{n}\}\) 的前 \(2n\) 项和 \(S_{2n}\).
\end{enumerate}
\end{problem}

\begin{solution}
\begin{enumerate}
\item 对任意正整数 \(k\)，有
\[
x^2-(3k+2^k)x+3k\cdot 2^k=(x-3k)(x-2^k)
\]
所以 \(a_{2k-1}\)，\(a_{2k}\) 是 \(3k\)，\(2^k\) 中较小者和较大者.
当 \(k=1\) 时，\(a_1=2\)，\(a_2=3\)；当 \(k=2\) 时，\(a_3=4\)，\(a_4=6\)；当 \(k=3\) 时，\(a_5=8\)，\(a_6=9\)；当 \(k=4\) 时，\(a_7=12\)，\(a_8=16\).
因此 \(a_1=2\)，\(a_3=4\)，\(a_5=8\)，\(a_7=12\).
当 \(n=4\) 时，\(2^n=16>12=3n\). 若 \(2^n>3n\)，则因 \(n\geqslant4\)，有 \(2^{n+1}>6n>3(n+1)\)，故对 \(n\geqslant4\) 恒有 \(2^n>3n\)，从而
\[
a_{2n}=2^n
\]

\item 每一对相邻项的和等于方程两根之和，因此
\[
S_{2n}=\sum_{k=1}^{n}(a_{2k-1}+a_{2k})=\sum_{k=1}^{n}(3k+2^k)
\]
又
\[
S_{2n}=3\cdot\frac{n(n+1)}{2}+2^{n+1}-2=\frac{3n(n+1)}{2}+2^{n+1}-2
\]
\end{enumerate}
\end{solution}

\begin{problem}
在如图所示的几何体中，\(EA \perp\) 平面 \(ABC\)，\(DB \perp\) 平面 \(ABC\)，\(AC \perp BC\)，且 \(AC = BC = BD = 2AE\)，\(M\) 是 \(AB\) 的中点.
\begin{enumerate}
\item 求证 \(CM \perp EM\)；
\item 求 \(DE\) 与平面 \(EMC\) 所成角的正切值.
\end{enumerate}

\bitmapfigure[max width=6.0cm]{img/2007/zhejiang_liberal/q20_fig1.png}
\end{problem}

\begin{solution}
\begin{enumerate}
\item 建立如图所示的直角坐标系，使平面 \(ABC\) 为 \(z=0\)，并取 \(C(0,0,0)\)，\(A(2,0,0)\)，\(B(0,2,0)\).
由 \(AC=BC=BD=2AE\)，可令 \(AE=1\)，于是
\(E(2,0,1)\)，\(D(0,2,2)\)，\(M(1,1,0)\).
因此 \(\overrightarrow{CM}=(1,1,0)\)，\(\overrightarrow{EM}=(-1,1,-1)\)，
从而
\[
\overrightarrow{CM}\cdot\overrightarrow{EM}=1\times(-1)+1\times1+0\times(-1)=0
\]
所以 \(CM\perp EM\).

\item 设直线 \(DE\) 与平面 \(EMC\) 所成的角为 \(\theta\). 取直线 \(DE\) 的方向向量以及平面 \(EMC\) 的一个法向量分别为
\(\bs{d}=\overrightarrow{DE}=(2,-2,-1)\)，\(\bs{n}=\overrightarrow{CM}\times\overrightarrow{EM}=(-1,1,2)\).
于是
\(|\bs{d}|=3\)，\(|\bs{n}|=\sqrt{6}\)，\(|\bs{d}\cdot\bs{n}|=6\).
由直线与平面所成角的定义，
\[
\sin\theta=\frac{|\bs{d}\cdot\bs{n}|}{|\bs{d}|\,|\bs{n}|}=\frac{6}{3\sqrt{6}}=\frac{\sqrt{6}}{3}
\]
因此 \(\cos\theta=\sqrt{1-\frac{6}{9}}=\frac{1}{\sqrt{3}}\)，故
\[
\tan\theta=\frac{\sin\theta}{\cos\theta}=\sqrt{2}
\]
\end{enumerate}
\end{solution}

\begin{problem}
如图，直线 \( y = kx + b \) 与椭圆 \( \frac{x^2}{4} + y^2 = 1 \) 交于 \(A\)，\(B\) 两点，记 \( \triangle AOB \) 的面积为 \(S\).
\begin{enumerate}
\item 求在 \(k = 0\)，\(0 < b < 1\) 的条件下，\(S\) 的最大值；
\item 当 \(\left|AB\right|=2\)，\(S=1\) 时，求直线 \(AB\) 的方程.
\end{enumerate}

\bitmapfigure[max width=6.2cm]{img/2007/zhejiang_liberal/q21_fig1.png}
\end{problem}

\begin{solution}
\begin{enumerate}
\item 当 \(k=0\) 时，直线为 \(y=b\). 由椭圆方程得交点的横坐标为 \(x=\pm2\sqrt{1-b^2}\)，所以
\[
|AB|=4\sqrt{1-b^2}
\]
三角形 \(AOB\) 的高为 \(b\)，故
\[
S=\frac{1}{2}\cdot4\sqrt{1-b^2}\cdot b=2b\sqrt{1-b^2}
\]
因为 \(0<b<1\)，有
\[
S^2=4b^2(1-b^2)\leqslant1
\]
当且仅当 \(b^2=1-b^2\)，即 \(b=\frac{\sqrt{2}}{2}\) 时取等号. 所以 \(S\) 的最大值为 \(1\).

\item 直线 \(y=kx+b\) 到原点的距离为 \(\dfrac{|b|}{\sqrt{1+k^2}}\). 由 \(|AB|=2\) 和 \(S=1\)，得
\[
\frac{1}{2}|AB|\cdot\frac{|b|}{\sqrt{1+k^2}}=1
\]
即
\[
b^2=k^2+1
\]
将直线方程代入椭圆方程，得交点横坐标满足
\[
\left(k^2+\frac{1}{4}\right)x^2+2kbx+b^2-1=0
\]
设此方程的两根为 \(x_1\)，\(x_2\)，则
\[
|AB|=|x_1-x_2|\sqrt{1+k^2}
\]
且其判别式为
\[
\Delta=(2kb)^2-4\left(k^2+\frac{1}{4}\right)(b^2-1)=4k^2+1-b^2
\]
所以
\[
|AB|=\frac{\sqrt{4k^2+1-b^2}\sqrt{1+k^2}}{k^2+\frac{1}{4}}
\]
代入 \(b^2=k^2+1\) 及 \(|AB|=2\)，得
\[
2=\frac{\sqrt{3k^2}\sqrt{1+k^2}}{k^2+\frac{1}{4}}
\]
两边平方并整理，得
\(3k^2(1+k^2)=4\left(k^2+\frac{1}{4}\right)^2\)，\(\left(k^2-\frac{1}{2}\right)^2=0\).
因此 \(k=\pm\frac{\sqrt{2}}{2}\)，且 \(b=\pm\sqrt{\frac{3}{2}}=\pm\frac{\sqrt{6}}{2}\). 这些取值均满足原条件，故直线 \(AB\) 为
\[
y=\pm\frac{\sqrt{2}}{2}x\pm\frac{\sqrt{6}}{2}
\]
其中两个符号相互独立.
\end{enumerate}
\end{solution}

\begin{problem}
已知 \( f(x) = |x^2 - 1| + x^2 + kx \).

\begin{enumerate}
\item 若 \(k = 2\)，求方程 \(f(x) = 0\) 的解；
\item 若关于 \(x\) 的方程 \(f(x) = 0\) 在 \((0,2)\) 上有两个解 \(x_{1}\)，\(x_{2}\)，求 \(k\) 的取值范围，并证明 \(\frac{1}{x_1} + \frac{1}{x_2} < 4\).
\end{enumerate}
\end{problem}

\begin{solution}
\begin{enumerate}
\item 由绝对值的定义，
\[
f(x)=
\begin{cases}
2x^2+kx-1,& |x|\geqslant1,\\
1+kx,& |x|<1.
\end{cases}
\]
当 \(k=2\) 时，在 \(|x|\geqslant1\) 上，方程 \(f(x)=0\) 化为
\[
2x^2+2x-1=0
\]
解得 \(x=\dfrac{-1\pm\sqrt{3}}{2}\). 其中 \(\dfrac{-1-\sqrt{3}}{2}\leqslant-1\)，而 \(\dfrac{-1+\sqrt{3}}{2}\in(0,1)\)，故前者符合本段定义域.
在 \(|x|<1\) 上，方程化为 \(1+2x=0\)，解得 \(x=-\frac{1}{2}\)，符合本段定义域. 所以方程的解为
\[
x=-\frac{1+\sqrt{3}}{2}\quad\text{或}\quad x=-\frac{1}{2}
\]

\item 记
\[
q(x)=2x^2+kx-1
\]
在区间 \((0,2)\) 中还需检查边界点 \(x=1\). 若 \(x=1\) 是方程的解，则 \(f(1)=1+k=0\)，即 \(k=-1\). 此时在 \(0<x<1\) 时 \(f(x)=1-x>0\)，在 \(1<x<2\) 时 \(f(x)=2x^2-x-1=(2x+1)(x-1)>0\)，所以 \((0,2)\) 中只有一个解 \(x=1\)，不能满足有两个解.
在 \(0<x<1\) 时，方程 \(f(x)=0\) 为 \(1+kx=0\)，其解为 \(x=-\dfrac{1}{k}\). 该解属于 \((0,1)\) 的充要条件是 \(k<-1\).

在 \(1<x<2\) 时，方程为 \(q(x)=0\). 因为 \(q(0)=-1<0\)，且二次方程 \(q(x)=0\) 的两根之积为 \(-\frac{1}{2}<0\)，所以它恰有一个正根. 这个正根属于 \((1,2)\) 的充要条件是
\(q(1)=k+1<0\)，\(q(2)=2k+7>0\)
即 \(-\frac{7}{2}<k<-1\). 在这个范围内，\(0<x<1\) 和 \(1<x<2\) 各有一个解，所以方程在 \((0,2)\) 上有两个解；反之若有两个解，也必须分别来自这两个区间，故
\[
-\frac{7}{2}<k<-1
\]

不妨设 \(x_1\in(0,1)\)，\(x_2\in(1,2)\). 由前面的方程，
\(x_1=-\frac{1}{k}\)，\(2x_2^2+kx_2-1=0\).
由第二式除以 \(x_2>0\)，得 \(\dfrac{1}{x_2}=k+2x_2\). 因而
\[
\frac{1}{x_1}+\frac{1}{x_2}=-k+k+2x_2=2x_2<4
\]
结论得证.
\end{enumerate}
\end{solution}
